\section{Einleitung}
In diesem Versuch wird die Wechselwirkung von $\alpha$-Strahlung mit Materie untersucht, insbesondere die Compton-Streuung. Ziel ist die Bestimmung des gesamten Wirkungsquerschnitts der Compton-Streuung und die Kalibrierung des verwendeten Spektrometers. Anschließend werden die Streuspektren einer $^{137}$Cs-Quelle aufgenommen und analysiert. Der Schwerpunkt des Experiments liegt auf der detaillierten Datenanalyse.


Zunächst wird das Ausgangssignal an der Anode des Szintillationsdetektors auf dem Oszilloskop beobachtet. Anschließend wird der Wirkungsquerschnitt bestimmt, indem der Szintillationsdetektor parallel zur $\gamma$-Quelle ($^{137}$Cs) hinter dem Bleikollimator positioniert und das Impulshöhenspektrum für Aluminiumtargets mit unterschiedlichen Durchmessern aufgezeichnet wird.

Zusätzlich wird eine Energiekalibrierung durchgeführt, wobei der Szintillationsdetektor senkrecht zur $^{137}$Cs-Quelle ausgerichtet ist. Hierzu wird eine $^{133}$Ba-Probe vor dem Detektor platziert und Messungen sowohl mit als auch ohne Probe durchgeführt.

Im letzten Teil des Experiments wird die $^{137}$Cs-Quelle erneut verwendet, um das Impulshöhenspektrum für verschiedene Streuwinkel, sowohl mit als auch ohne das runde Target, aufzuzeichnen.